% 核心结论：势能的负梯度给力，电势的负梯度给电场。
% 图契约：物理竞赛讲义；印刷宽16cm，最小字号8.5pt；几何示意与公式分区。
% 布局：[几何/曲线 7cm] [定义/结果 7cm]；实线主量，虚线辅助；无装饰数据。
\documentclass[tikz,border=6pt]{standalone}
\usepackage{fontspec,xeCJK,amsmath,amssymb,graphicx}
\IfFontExistsTF{Times New Roman}{\setmainfont{Times New Roman}}{\setmainfont{TeX Gyre Termes}}
\IfFontExistsTF{Heiti SC}{\setCJKmainfont[AutoFakeBold=1.5]{Heiti SC}}{\setCJKmainfont[AutoFakeBold=1.5]{FandolHei-Regular}}
\usetikzlibrary{arrows.meta,bending,calc,positioning,patterns}
\definecolor{ink}{HTML}{243745}
\definecolor{blue}{HTML}{275DAB}
\definecolor{green}{HTML}{177C63}
\definecolor{red}{HTML}{BD4935}
\definecolor{orange}{HTML}{AA6C16}
\definecolor{purple}{HTML}{715393}
\tikzset{
 every node/.style={font=\fontsize{9.5}{13}\selectfont,text=ink},
 axis/.style={-{Stealth[length=4pt]},line width=.65pt,draw=ink!65},
 vec/.style={-{Stealth[length=5pt]},line width=1pt,draw=blue},
 short/.style={-{Stealth[length=3pt]},line width=.8pt},
 guide/.style={dashed,draw=ink!40,line width=.5pt},
 card/.style={rounded corners=2pt,draw=blue!60,fill=blue!4,
   align=center,text width=6.8cm,inner xsep=9pt,inner ysep=7pt},
 title/.style={font=\fontsize{11}{14}\selectfont\bfseries,text width=15cm,align=center},
 label/.style={font=\fontsize{8.5}{11}\selectfont,align=center}}
\begin{document}
\begin{tikzpicture}[x=1cm,y=1cm]
\path[use as bounding box] (-8,-3.35) rectangle (8,3.35);
\node[title] at (0,2.95) {负梯度：区分势能、电势、电场与力};

\node[card,text width=5.7cm] (u) at (-4.1,1.55) {给定粒子的势能 $U$\\$\mathbf F=-\nabla U$};
\node[card,text width=5.7cm,draw=green!65,fill=green!4] (p) at (4.1,1.55) {静电场的电势 $\varphi$\\$\mathbf E=-\nabla\varphi$};
\node[card,text width=13.5cm] at (0,-.1) {$U=q\varphi$，$\mathbf F=q\mathbf E$\\$q>0$ 时力与电场同向；$q<0$ 时力与电场反向。};
\draw[vec] (-5,-1.75)--(4.8,-1.75) node[right] {$\mathbf E$};
\node[label] at (0,-2.45) {沿电场方向电势降低；势能的变化还取决于试探电荷的符号。};

\end{tikzpicture}\end{document}
